\documentclass[a4paper,12pt]{article}
\usepackage{amsmath, amssymb}

\title{Løsning av $Ax = b$ med radoperasjoner}
\author{}
\date{}

\begin{document}

\maketitle

\section*{Gitt}

La matrisen $A$ og vektoren $b$ være:

\[
A =
\begin{bmatrix}
-3 & 0 & 0 \\
3 & 2 & -3 \\
2 & -2 & 3
\end{bmatrix},
\quad
b =
\begin{bmatrix}
b_1 \\
b_2 \\
b_3
\end{bmatrix}
\]

Vi ønsker å løse ligningssystemet $Ax = b$ ved hjelp av radoperasjoner (Gauss-eliminasjon).

\section*{Steg 1: Utvidet matrise $[A \mid b]$}

\[
\left[
\begin{array}{ccc|c}
-3 & 0 & 0 & b_1 \\
3 & 2 & -3 & b_2 \\
2 & -2 & 3 & b_3
\end{array}
\right]
\]

\section*{Steg 2: Gjør første pivot lik 1}

Rad 1 divideres med $-3$:

\[
R_1 \rightarrow \frac{1}{-3} R_1
\Rightarrow
\left[
\begin{array}{ccc|c}
1 & 0 & 0 & -\frac{b_1}{3} \\
3 & 2 & -3 & b_2 \\
2 & -2 & 3 & b_3
\end{array}
\right]
\]

\section*{Steg 3: Null ut under pivot i kolonne 1}

\begin{align*}
R_2 &\rightarrow R_2 - 3R_1 \\
R_3 &\rightarrow R_3 - 2R_1
\end{align*}

\[
\left[
\begin{array}{ccc|c}
1 & 0 & 0 & -\frac{b_1}{3} \\
0 & 2 & -3 & b_2 + b_1 \\
0 & -2 & 3 & b_3 + \frac{2b_1}{3}
\end{array}
\right]
\]

\section*{Steg 4: Gjør pivot i R2 til 1}

\[
R_2 \rightarrow \frac{1}{2}R_2
\Rightarrow
\left[
\begin{array}{ccc|c}
1 & 0 & 0 & -\frac{b_1}{3} \\
0 & 1 & -\frac{3}{2} & \frac{b_2 + b_1}{2} \\
0 & -2 & 3 & b_3 + \frac{2b_1}{3}
\end{array}
\right]
\]

\section*{Steg 5: Null ut under pivot i kolonne 2}

\[
R_3 \rightarrow R_3 + 2R_2
\Rightarrow
\left[
\begin{array}{ccc|c}
1 & 0 & 0 & -\frac{b_1}{3} \\
0 & 1 & -\frac{3}{2} & \frac{b_2 + b_1}{2} \\
0 & 0 & 0 & b_2 + b_3 + \frac{5b_1}{3}
\end{array}
\right]
\]

\section*{Tolkning}

Den siste raden gir:
\[
0 = b_2 + b_3 + \frac{5b_1}{3}
\]

Systemet er konsistent (har løsning) hvis og bare hvis:
\[
b_2 + b_3 = -\frac{5b_1}{3}
\]

\section*{Generell løsning (når systemet er løsbart)}

Anta $z = t$, en fri parameter.

\begin{align*}
x &= -\frac{b_1}{3} \\
y &= \frac{b_2 + b_1}{2} + \frac{3}{2}t \\
z &= t
\end{align*}

\section*{Løsningsmengde}

Løsningsmengden kan skrives som:

\[
\begin{bmatrix}
x \\
y \\
z
\end{bmatrix}
=
\begin{bmatrix}
-\frac{b_1}{3} \\
\frac{b_2 + b_1}{2} \\
0
\end{bmatrix}
+ t \cdot
\begin{bmatrix}
0 \\
\frac{3}{2} \\
1
\end{bmatrix}, \quad t \in \mathbb{R}
\]

\section*{Oppsummering}

\begin{itemize}
    \item Systemet har løsning hvis og bare hvis: \( b_2 + b_3 = -\frac{5b_1}{3} \)
    \item Da finnes uendelig mange løsninger med én fri parameter \( t \in \mathbb{R} \)
    \item Løsningen kan uttrykkes som en vektor med parameterframstilling
\end{itemize}

\end{document}